% Evaluation methodology — owner: Abhayanth.
% Source of truth: docs/EVALUATION_PROTOCOL.md. Code: src/rul_prediction/{targets,splits,metrics}.py.
% Include from the main report with % Evaluation methodology — owner: Abhayanth.
% Source of truth: docs/EVALUATION_PROTOCOL.md. Code: src/rul_prediction/{targets,splits,metrics}.py.
% Include from the main report with % Evaluation methodology — owner: Abhayanth.
% Source of truth: docs/EVALUATION_PROTOCOL.md. Code: src/rul_prediction/{targets,splits,metrics}.py.
% Include from the main report with % Evaluation methodology — owner: Abhayanth.
% Source of truth: docs/EVALUATION_PROTOCOL.md. Code: src/rul_prediction/{targets,splits,metrics}.py.
% Include from the main report with \input{sections/evaluation_method}.
% Needs: amsmath.

\section{Evaluation Methodology}
\label{sec:evaluation}

All models in this study (linear, tree-based and sequence models) share one
target definition, one engine-level data split, one set of validation points
and one pair of metric functions, implemented once in shared code. This makes
their errors directly comparable.

\subsection{Target Definition}

Each training engine in FD001 is run until failure, so its last recorded cycle
is its failure point. For engine $i$ with final cycle $T_i$, the Remaining
Useful Life (RUL) at cycle $t$ is
\begin{equation}
  \mathrm{RUL}_{i,t} = T_i - t ,
  \label{eq:rul}
\end{equation}
which falls to 0 at the failure cycle. Training trajectories last between 128
and 362 cycles, so $T_i$ is looked up per engine. No upper cap is applied to
the label. A cap (for example 125 cycles) changes the learning target, so it
is treated as a separate experiment rather than a default. Both the training
target and the answer key stay true RUL. A capped training target would
replace the default only if its overall validation RMSE against true RUL were
no worse than without the cap and its late-band error (defined below) were
lower. This rule was fixed before the cap experiment was run.

Test engines stop before failure. Only their last observed cycle has a known
RUL, supplied in \texttt{RUL\_FD001.txt}. Each test engine is therefore scored
once, at its last observed cycle. Predictions are joined to the answer key by
engine ID rather than by row position, and the join fails if any engine is
missing, extra or repeated.

\subsection{Engine-Level Validation Split}

Consecutive cycles of one engine have almost identical sensor readings. A
random row split would place the same engine in both training and validation,
letting a model memorise that engine and report an optimistic score. We
therefore split by engine: 20\% of the 100 training engines (20 engines,
random seed 42) are held out for validation. This leaves 16{,}390 training
rows and 4{,}241 validation rows, with no engine on both sides. Scalers,
feature selection and any other fitted preprocessing are fitted on the 80
training engines only.

\subsection{Test-Like Validation Points}

Scoring every validation row would not resemble the test set, where each
engine is seen only up to an unknown cut-off before failure. Instead, for each
validation engine we sample 5 cut-off cycles uniformly at random (seed 42)
between cycle 31 and $T_i - 1$, which gives 100 validation points. The lower
bound of 31 is the shortest test-engine history; it is derived from test
input lengths only, never from test labels. The upper bound guarantees that
the engine has not yet failed.

Scoring the row at cut-off $c$ is equivalent to truncating the engine at $c$
only if every feature is causal, that is, computed from cycles $\le c$ of the
same engine. Centred or forward-looking windows and whole-trajectory
statistics are therefore not allowed. Sequence models use the window that
ends at each cut-off. The 100 points are committed to the repository and
pinned by an automated test, so every model is scored on identical rows.

\subsection{Metrics}

For $n$ scored points with true RUL $y_k$ and prediction $\hat{y}_k$:
\begin{align}
  \mathrm{MAE}  &= \frac{1}{n}\sum_{k=1}^{n} \lvert y_k - \hat{y}_k \rvert , \\
  \mathrm{RMSE} &= \sqrt{\frac{1}{n}\sum_{k=1}^{n} \left(y_k - \hat{y}_k\right)^2 } .
\end{align}
Both are measured in cycles, and lower is better. MAE is the average size of
a miss. RMSE squares errors before averaging, so a few large misses raise it
more than many small ones; $\mathrm{RMSE} \ge \mathrm{MAE}$ always holds. Over-
and under-prediction are penalised equally. Hyperparameters are selected by
validation RMSE, because large misses are the costly ones in maintenance
planning.

To show where the errors come from, each set of predictions is also scored in
two bands, grouped by the true RUL after prediction: \emph{early} points
(true RUL $> 125$, engine assumed still healthy) and \emph{late} points (true
RUL $\le 125$, engine assumed degrading). The threshold is borrowed from the
capped targets common in C-MAPSS work, which go back to the piecewise-linear
target of Heimes (2008). It is not a measured onset of degradation, and it was
fixed before any band score was computed. The 100 validation points split
into 22 early points from 14 engines and 78 late points from all 20 engines.
The overall MAE and RMSE remain the primary scores. Test-set band scores are
computed only in the final evaluation, after the model list is frozen.

\subsection{Use of the Test Set}

Models, features and hyperparameters are chosen on the validation points
only. Each final model is then evaluated on the test set once, and it is not
changed after its test score is seen. Test RUL values are not summarised or
plotted before that final evaluation.

\subsection{Limitations of the Protocol}

The 100 validation points come from only 20 engines, and points from the same
engine are correlated. The effective sample size is therefore closer to 20,
and differences of a few cycles between models may be noise. Tuning
repeatedly against one fixed split can also make validation scores
optimistic. Finally, uniform cut-offs include early-life points with RUL up to
about 330 cycles, which weigh heavily in uncapped RMSE.
.
% Needs: amsmath.

\section{Evaluation Methodology}
\label{sec:evaluation}

All models in this study (linear, tree-based and sequence models) share one
target definition, one engine-level data split, one set of validation points
and one pair of metric functions, implemented once in shared code. This makes
their errors directly comparable.

\subsection{Target Definition}

Each training engine in FD001 is run until failure, so its last recorded cycle
is its failure point. For engine $i$ with final cycle $T_i$, the Remaining
Useful Life (RUL) at cycle $t$ is
\begin{equation}
  \mathrm{RUL}_{i,t} = T_i - t ,
  \label{eq:rul}
\end{equation}
which falls to 0 at the failure cycle. Training trajectories last between 128
and 362 cycles, so $T_i$ is looked up per engine. No upper cap is applied to
the label. A cap (for example 125 cycles) changes the learning target, so it
is treated as a separate experiment rather than a default. Both the training
target and the answer key stay true RUL. A capped training target would
replace the default only if its overall validation RMSE against true RUL were
no worse than without the cap and its late-band error (defined below) were
lower. This rule was fixed before the cap experiment was run.

Test engines stop before failure. Only their last observed cycle has a known
RUL, supplied in \texttt{RUL\_FD001.txt}. Each test engine is therefore scored
once, at its last observed cycle. Predictions are joined to the answer key by
engine ID rather than by row position, and the join fails if any engine is
missing, extra or repeated.

\subsection{Engine-Level Validation Split}

Consecutive cycles of one engine have almost identical sensor readings. A
random row split would place the same engine in both training and validation,
letting a model memorise that engine and report an optimistic score. We
therefore split by engine: 20\% of the 100 training engines (20 engines,
random seed 42) are held out for validation. This leaves 16{,}390 training
rows and 4{,}241 validation rows, with no engine on both sides. Scalers,
feature selection and any other fitted preprocessing are fitted on the 80
training engines only.

\subsection{Test-Like Validation Points}

Scoring every validation row would not resemble the test set, where each
engine is seen only up to an unknown cut-off before failure. Instead, for each
validation engine we sample 5 cut-off cycles uniformly at random (seed 42)
between cycle 31 and $T_i - 1$, which gives 100 validation points. The lower
bound of 31 is the shortest test-engine history; it is derived from test
input lengths only, never from test labels. The upper bound guarantees that
the engine has not yet failed.

Scoring the row at cut-off $c$ is equivalent to truncating the engine at $c$
only if every feature is causal, that is, computed from cycles $\le c$ of the
same engine. Centred or forward-looking windows and whole-trajectory
statistics are therefore not allowed. Sequence models use the window that
ends at each cut-off. The 100 points are committed to the repository and
pinned by an automated test, so every model is scored on identical rows.

\subsection{Metrics}

For $n$ scored points with true RUL $y_k$ and prediction $\hat{y}_k$:
\begin{align}
  \mathrm{MAE}  &= \frac{1}{n}\sum_{k=1}^{n} \lvert y_k - \hat{y}_k \rvert , \\
  \mathrm{RMSE} &= \sqrt{\frac{1}{n}\sum_{k=1}^{n} \left(y_k - \hat{y}_k\right)^2 } .
\end{align}
Both are measured in cycles, and lower is better. MAE is the average size of
a miss. RMSE squares errors before averaging, so a few large misses raise it
more than many small ones; $\mathrm{RMSE} \ge \mathrm{MAE}$ always holds. Over-
and under-prediction are penalised equally. Hyperparameters are selected by
validation RMSE, because large misses are the costly ones in maintenance
planning.

To show where the errors come from, each set of predictions is also scored in
two bands, grouped by the true RUL after prediction: \emph{early} points
(true RUL $> 125$, engine assumed still healthy) and \emph{late} points (true
RUL $\le 125$, engine assumed degrading). The threshold is borrowed from the
capped targets common in C-MAPSS work, which go back to the piecewise-linear
target of Heimes (2008). It is not a measured onset of degradation, and it was
fixed before any band score was computed. The 100 validation points split
into 22 early points from 14 engines and 78 late points from all 20 engines.
The overall MAE and RMSE remain the primary scores. Test-set band scores are
computed only in the final evaluation, after the model list is frozen.

\subsection{Use of the Test Set}

Models, features and hyperparameters are chosen on the validation points
only. Each final model is then evaluated on the test set once, and it is not
changed after its test score is seen. Test RUL values are not summarised or
plotted before that final evaluation.

\subsection{Limitations of the Protocol}

The 100 validation points come from only 20 engines, and points from the same
engine are correlated. The effective sample size is therefore closer to 20,
and differences of a few cycles between models may be noise. Tuning
repeatedly against one fixed split can also make validation scores
optimistic. Finally, uniform cut-offs include early-life points with RUL up to
about 330 cycles, which weigh heavily in uncapped RMSE.
.
% Needs: amsmath.

\section{Evaluation Methodology}
\label{sec:evaluation}

All models in this study (linear, tree-based and sequence models) share one
target definition, one engine-level data split, one set of validation points
and one pair of metric functions, implemented once in shared code. This makes
their errors directly comparable.

\subsection{Target Definition}

Each training engine in FD001 is run until failure, so its last recorded cycle
is its failure point. For engine $i$ with final cycle $T_i$, the Remaining
Useful Life (RUL) at cycle $t$ is
\begin{equation}
  \mathrm{RUL}_{i,t} = T_i - t ,
  \label{eq:rul}
\end{equation}
which falls to 0 at the failure cycle. Training trajectories last between 128
and 362 cycles, so $T_i$ is looked up per engine. No upper cap is applied to
the label. A cap (for example 125 cycles) changes the learning target, so it
is treated as a separate experiment rather than a default. Both the training
target and the answer key stay true RUL. A capped training target would
replace the default only if its overall validation RMSE against true RUL were
no worse than without the cap and its late-band error (defined below) were
lower. This rule was fixed before the cap experiment was run.

Test engines stop before failure. Only their last observed cycle has a known
RUL, supplied in \texttt{RUL\_FD001.txt}. Each test engine is therefore scored
once, at its last observed cycle. Predictions are joined to the answer key by
engine ID rather than by row position, and the join fails if any engine is
missing, extra or repeated.

\subsection{Engine-Level Validation Split}

Consecutive cycles of one engine have almost identical sensor readings. A
random row split would place the same engine in both training and validation,
letting a model memorise that engine and report an optimistic score. We
therefore split by engine: 20\% of the 100 training engines (20 engines,
random seed 42) are held out for validation. This leaves 16{,}390 training
rows and 4{,}241 validation rows, with no engine on both sides. Scalers,
feature selection and any other fitted preprocessing are fitted on the 80
training engines only.

\subsection{Test-Like Validation Points}

Scoring every validation row would not resemble the test set, where each
engine is seen only up to an unknown cut-off before failure. Instead, for each
validation engine we sample 5 cut-off cycles uniformly at random (seed 42)
between cycle 31 and $T_i - 1$, which gives 100 validation points. The lower
bound of 31 is the shortest test-engine history; it is derived from test
input lengths only, never from test labels. The upper bound guarantees that
the engine has not yet failed.

Scoring the row at cut-off $c$ is equivalent to truncating the engine at $c$
only if every feature is causal, that is, computed from cycles $\le c$ of the
same engine. Centred or forward-looking windows and whole-trajectory
statistics are therefore not allowed. Sequence models use the window that
ends at each cut-off. The 100 points are committed to the repository and
pinned by an automated test, so every model is scored on identical rows.

\subsection{Metrics}

For $n$ scored points with true RUL $y_k$ and prediction $\hat{y}_k$:
\begin{align}
  \mathrm{MAE}  &= \frac{1}{n}\sum_{k=1}^{n} \lvert y_k - \hat{y}_k \rvert , \\
  \mathrm{RMSE} &= \sqrt{\frac{1}{n}\sum_{k=1}^{n} \left(y_k - \hat{y}_k\right)^2 } .
\end{align}
Both are measured in cycles, and lower is better. MAE is the average size of
a miss. RMSE squares errors before averaging, so a few large misses raise it
more than many small ones; $\mathrm{RMSE} \ge \mathrm{MAE}$ always holds. Over-
and under-prediction are penalised equally. Hyperparameters are selected by
validation RMSE, because large misses are the costly ones in maintenance
planning.

To show where the errors come from, each set of predictions is also scored in
two bands, grouped by the true RUL after prediction: \emph{early} points
(true RUL $> 125$, engine assumed still healthy) and \emph{late} points (true
RUL $\le 125$, engine assumed degrading). The threshold is borrowed from the
capped targets common in C-MAPSS work, which go back to the piecewise-linear
target of Heimes (2008). It is not a measured onset of degradation, and it was
fixed before any band score was computed. The 100 validation points split
into 22 early points from 14 engines and 78 late points from all 20 engines.
The overall MAE and RMSE remain the primary scores. Test-set band scores are
computed only in the final evaluation, after the model list is frozen.

\subsection{Use of the Test Set}

Models, features and hyperparameters are chosen on the validation points
only. Each final model is then evaluated on the test set once, and it is not
changed after its test score is seen. Test RUL values are not summarised or
plotted before that final evaluation.

\subsection{Limitations of the Protocol}

The 100 validation points come from only 20 engines, and points from the same
engine are correlated. The effective sample size is therefore closer to 20,
and differences of a few cycles between models may be noise. Tuning
repeatedly against one fixed split can also make validation scores
optimistic. Finally, uniform cut-offs include early-life points with RUL up to
about 330 cycles, which weigh heavily in uncapped RMSE.
.
% Needs: amsmath.

\section{Evaluation Methodology}
\label{sec:evaluation}

All models in this study (linear, tree-based and sequence models) share one
target definition, one engine-level data split, one set of validation points
and one pair of metric functions, implemented once in shared code. This makes
their errors directly comparable.

\subsection{Target Definition}

Each training engine in FD001 is run until failure, so its last recorded cycle
is its failure point. For engine $i$ with final cycle $T_i$, the Remaining
Useful Life (RUL) at cycle $t$ is
\begin{equation}
  \mathrm{RUL}_{i,t} = T_i - t ,
  \label{eq:rul}
\end{equation}
which falls to 0 at the failure cycle. Training trajectories last between 128
and 362 cycles, so $T_i$ is looked up per engine. No upper cap is applied to
the label. A cap (for example 125 cycles) changes the learning target, so it
is treated as a separate experiment rather than a default. Both the training
target and the answer key stay true RUL. A capped training target would
replace the default only if its overall validation RMSE against true RUL were
no worse than without the cap and its late-band error (defined below) were
lower. This rule was fixed before the cap experiment was run.

Test engines stop before failure. Only their last observed cycle has a known
RUL, supplied in \texttt{RUL\_FD001.txt}. Each test engine is therefore scored
once, at its last observed cycle. Predictions are joined to the answer key by
engine ID rather than by row position, and the join fails if any engine is
missing, extra or repeated.

\subsection{Engine-Level Validation Split}

Consecutive cycles of one engine have almost identical sensor readings. A
random row split would place the same engine in both training and validation,
letting a model memorise that engine and report an optimistic score. We
therefore split by engine: 20\% of the 100 training engines (20 engines,
random seed 42) are held out for validation. This leaves 16{,}390 training
rows and 4{,}241 validation rows, with no engine on both sides. Scalers,
feature selection and any other fitted preprocessing are fitted on the 80
training engines only.

\subsection{Test-Like Validation Points}

Scoring every validation row would not resemble the test set, where each
engine is seen only up to an unknown cut-off before failure. Instead, for each
validation engine we sample 5 cut-off cycles uniformly at random (seed 42)
between cycle 31 and $T_i - 1$, which gives 100 validation points. The lower
bound of 31 is the shortest test-engine history; it is derived from test
input lengths only, never from test labels. The upper bound guarantees that
the engine has not yet failed.

Scoring the row at cut-off $c$ is equivalent to truncating the engine at $c$
only if every feature is causal, that is, computed from cycles $\le c$ of the
same engine. Centred or forward-looking windows and whole-trajectory
statistics are therefore not allowed. Sequence models use the window that
ends at each cut-off. The 100 points are committed to the repository and
pinned by an automated test, so every model is scored on identical rows.

\subsection{Metrics}

For $n$ scored points with true RUL $y_k$ and prediction $\hat{y}_k$:
\begin{align}
  \mathrm{MAE}  &= \frac{1}{n}\sum_{k=1}^{n} \lvert y_k - \hat{y}_k \rvert , \\
  \mathrm{RMSE} &= \sqrt{\frac{1}{n}\sum_{k=1}^{n} \left(y_k - \hat{y}_k\right)^2 } .
\end{align}
Both are measured in cycles, and lower is better. MAE is the average size of
a miss. RMSE squares errors before averaging, so a few large misses raise it
more than many small ones; $\mathrm{RMSE} \ge \mathrm{MAE}$ always holds. Over-
and under-prediction are penalised equally. Hyperparameters are selected by
validation RMSE, because large misses are the costly ones in maintenance
planning.

To show where the errors come from, each set of predictions is also scored in
two bands, grouped by the true RUL after prediction: \emph{early} points
(true RUL $> 125$, engine assumed still healthy) and \emph{late} points (true
RUL $\le 125$, engine assumed degrading). The threshold is borrowed from the
capped targets common in C-MAPSS work, which go back to the piecewise-linear
target of Heimes (2008). It is not a measured onset of degradation, and it was
fixed before any band score was computed. The 100 validation points split
into 22 early points from 14 engines and 78 late points from all 20 engines.
The overall MAE and RMSE remain the primary scores. Test-set band scores are
computed only in the final evaluation, after the model list is frozen.

\subsection{Use of the Test Set}

Models, features and hyperparameters are chosen on the validation points
only. Each final model is then evaluated on the test set once, and it is not
changed after its test score is seen. Test RUL values are not summarised or
plotted before that final evaluation.

\subsection{Limitations of the Protocol}

The 100 validation points come from only 20 engines, and points from the same
engine are correlated. The effective sample size is therefore closer to 20,
and differences of a few cycles between models may be noise. Tuning
repeatedly against one fixed split can also make validation scores
optimistic. Finally, uniform cut-offs include early-life points with RUL up to
about 330 cycles, which weigh heavily in uncapped RMSE.
